\documentclass{article}
\usepackage{amssymb}

\oddsidemargin 0 in
\evensidemargin 0 in
\textwidth 6.5 in
\topmargin 0 in
\textheight 8.5 in
\sloppy
\setlength{\parindent}{0.0in}
\setlength{\parskip}{7pt plus 1pt}

\begin{document}
\sloppy

\title{Object Identity Constraints in Babelsberg}
\author{Alan Borning and Tim Felgentreff}
\date{Draft of 10 October 2013}

\maketitle

\section{Introduction}

Babelsberg \cite{felgentreff-babelsberg-tr} is a family of object
constraint languages that integrate constraints, expressed in a
domain-specific language, with a standard host object-oriented
language.  Current versions are Babelsberg/R (Ruby) and Babelsberg/JS
(JavaScript/Lively).

Babelsberg supports constraints on object identity, as well as on
equality and other relations.  For example, this constraint ensures
that \verb|x| and \verb|y| refer to the same object.\footnote{Standard
  Ruby uses the method {\tt equal?} for the object identity test, but
  Babelsberg/R uses a new keyword, {\tt is?}, both to avoid
  programmers overriding the identity method, and also because
  programmers coming from other languages are likely to find Ruby's
  standard method name peculiar if not downright confusing.  {\tt ==}
  is Ruby's equality test.}  (In other words, there must be just one
object that both \verb|x| and \verb|y| refer to, not two equal
objects.)

\begin{verbatim}
x is? y
\end{verbatim}


However, there are more complicated examples of identity constraints
that either don't make sense semantically, or that do make sense but
are unsatisfiable for not completely obvious reasons.  This note tries
to clarify the rules for what is legal in a constraint involving
object identity.  

The Babelsberg paper \cite{felgentreff-babelsberg-tr} lists three
restrictions on constraints:

\begin{quotation}
First, evaluating the expression that defines the constraint should return a
boolean --- the constraint is that the expression evaluates to true.

Second, the constraint expression should either be free of side
effects, or if there are side effects, they should be benign (for
example, doing caching).  Also, repeatedly evaluating the block should
return the same thing (so an expression involving a random number
generator wouldn't qualify).  This restriction is needed to provide
the correct semantics for constraints.

Third, variables used in methods that will be called in a constraint
must be used in a single-assignment fashion.  
\end{quotation}

To handle identity constraints properly, the second restriction should
be strengthened, and an additional restriction should be added.  Here
is a proposed revision, which also divides the second restriction into
several parts and fixes a smaller issue regarding system primitives.

\begin{enumerate}

\item \label{item:constraint-type} Evaluating the expression that
  defines the constraint should return a boolean --- the constraint is
  that the expression evaluates to true.

\item \label{item:side-effects} The constraint expression should
  either be free of side effects, or if there are side effects, they
  should be benign (for example, doing caching).

\item \label{item:deterministic} The result of evaluating the block
  should be deterministic.  (For example, an expression whose value
  depended on which of two processes happened to complete first
  wouldn't qualify.)

\item \label{item:single-assignment} Variables used in methods that
  will be called in a constraint must be used in a single-assignment
  fashion.

\item \label{item:identity-constraints} If the block uses identity
  constraints, neither the constraint block nor any methods that it
  calls (directly or indirectly) can create new objects --- they can
  only use existing objects.

\end{enumerate}

Restrictions \ref{item:constraint-type}--\ref{item:deterministic} are
needed to provide the correct semantics for constraints.

Restriction \ref{item:single-assignment} (single assignment semantics)
prevents a class of errors resulting from overconstraining a
multiply-assigned-to variable.  See \cite{felgentreff-babelsberg-tr}
for more information, and alternatives.

Restriction \ref{item:identity-constraints} makes it easier for the
identity constraint solver to ensure that the identity constraints are
indeed satisfied.  An intuition for this restriction is that changing
the UID of the object bound to a variable is something like a side
effect if you have identity constraints.  Some examples are given in
Section \ref{sec:examples}.  It is actually a stronger restriction
than is needed, and we might be able to get rid of it altogether by
including runtime checks to make sure that the constraint really was
satisfied.  However, we don't have clear use cases for needing to
violate this restriction.  So we leave it in for now, so that we can
make progress, but may return to it later.  (See the discussion of
alternatives in Section \ref{sec:alternatives}.)

\section{Semantics of Constraints}

The semantics of constraints are straightforward: the expression, as
evaluable Ruby code, defines the constraint.  If it evaluates to true
the constraint is satisfied.

However, there are some less straightforward consequences with respect
to the solvers.  The Babelsberg/R VM, in concert with the solvers, is
charged with satisfying the constraints.  Given a required constraint
to solve, the solver should either satisfy it or raise an exception
(where ``satisfying it'' implies that the value resulting from
evaluating the expression should be true).  Ideally, the solver would
also include any runtime checks needed to ensure that it functioned
correctly for any legal constraint, but we might need to include
additional restrictions on the constraints so that the solver is known
to work with them.  (Considering soft constraints as well as hard or
required ones doesn't seem to add any additional issues, but since it
would complicate the discussion let's ignore them for now.)

\section{Implementation Considerations}
\label{sec:implementation}

The restriction that the result of evaluating the block be
deterministic implies that if the system is in the same state,
evaluating the constraint block should always give the same result.
In the Babelsberg/R implementation, operationally this means that the
value of the block should only depend on variables known to the
solver, and not on system methods that can return different values
without the solver being triggered.  This restriction could perhaps be
removed if the VM were suitably modified, for example so that the
solver was triggered if the value that would be returned by a
primitive method changed, but this is probably not worth the
performance penalty and implementation complexity.

In the current implementation, identity constraints in Babelsberg/R
are handled at the \textsc{vm} level rather than with a separate
solver.  An identity constraint adds links between the constrained
variables. An assignment to one of the variables is semantically the
same as a simultaneous assignment of the same value to all linked
variables.  (See the discussion on multi-assignments
in~\cite{felgentreff-babelsberg-tr} for more information on
simultaneous assignments.)  This simple propagation mechanism
currently does not allow constraint priorities on identity
constraints; however, this is only an implementation issue rather than
being a necessary aspect of the semantics.  Another limitation of the
current implementation (not the design) is that a constraint
expression can only contain identity constraints, or other constraints
such as equalities or inequalities, but not both.  

\section{Examples}
\label{sec:examples}

Let's look at some of the examples that have been discussed in emails
or draft memos.

The following code creates a new point \verb|p| and then requires it
to be equal (not identical) to a point with x=100, y=100:

\begin{verbatim}
p = Point.new(0,0)
always { p == Point.new(100,200) }
\end{verbatim}

In Babelsberg/R the interpreter can run in ordinary imperative
evaluation mode, or in constraint construction mode.  In ordinary
evaluation mode we can use the expression as a test, as discussed
above.  When in constraint construction mode, evaluating the
\verb|always| constraint will run point's == method (the equality
test), and add primitive constraints that p's x field equal 100, and
p's y field equal 200, keeping track of dependencies in the process.
This works fine, even though we are creating a new point each time we
use the constraint as a test.  In particular, if we repeatedly
evaluate the constraint block as a test, we always get true.

However, an otherwise similar constraint that uses an identity
constraint rather than an equality constraint is unsatisfiable:

\begin{verbatim}
p = Point.new(0,0)
always { p is? Point.new(100,200) }
\end{verbatim}

No matter what \verb|p| is bound to, running the constraint as a test
will return false, since we are testing whether \verb|p| is identical
to a freshly-created point.  Something potentially confusing about
this example is that we can imagine that the constraint solver would
create a new point with x=100 and y=200, happily bind it to p, and at
that time regard the constraint as satisfied.  But is the constraint
satisfied at that time?  If we then run the constraint block as a test
it returns false, so by our semantics, it is not.  The restrictions
listed above would catch this issue, since the constraint involves
creating a new object.  (As hinted above, one could probably also
catch this problem by first running the solver to satisfy the
constraint, then running the constraint as a test to make sure the
constraint really is satisfied.  But doing this seems unnecessary,
since so far we've only constructed pathological cases that need this
test.)

Next, consider this example from \cite{felgentreff-babelsberg-tr},
which uses just an equality constraint but not identity.  Suppose also
that rectangles are represented using origin/corner, so that
\verb|center| is a method.

\begin{verbatim}
always { rect.center == Point.new(100,100) }
\end{verbatim}

When in constraint construction mode, evaluating the \verb|always|
constraint will run rectangle's center method, and add primitive
constraints that \verb|(origin.x+corner.x)/2==100| and similarly for
\verb|y|, keeping track of dependencies in the process.  Using the
constraint as a test also works fine.

However, an otherwise similar constraint that uses an identity
constraint rather than an equality constraint is unsatisfiable, for
the same reasons as before, and disallowed by the prohibition against
creating new objects in constraints that use identity.

\begin{verbatim}
always { rect.center is? Point.new(100,100) }
\end{verbatim}

\subsection{Invoking Other Methods from Identity Constraints}

One or both sides of an identity constraint can invoke methods that
include conditionals, loops, and recursion --- they just can't create
new objects.  For example, we could have a method \verb|last| that
returns the last element of a linked list, either using a while loop
or recursive -- in either case, it is more than a simple accessor
method.  We could then use this method in a constraint like

\begin{verbatim}
always { e is? s.last }
\end{verbatim}

Identity constraints can use temporary variables or \verb|self|
(\verb|self| is treated as read-only --- rebinding \verb|self| to
satisfy an identity constraint would arguably be a side effect, and in
any case would be weird).


\subsection{Examples with System Calls}
\label{sec:system-calls}

Here are a few constraints that are disallowed because they call
system methods.  (They use equality rather than identity constraints
--- the aspect that is problematic is using primitive methods.)

\begin{verbatim} 
always { x == rand? }        # this uses the built-in rand method 
always { t == Time.now() }
always { b == File.exist?("fred.txt") }
\end{verbatim} 

\subsection{Stack Examples}

Suppose we have a class Stack with \verb|init|, \verb|push|,
\verb|pop|, and \verb|empty?| methods.

Then this constraint is disallowed because it would have a side
effect:

\begin{verbatim}
always { x == stack.pop }
\end{verbatim}

(And clearly it \emph{should} be disallowed --- every time we evaluate
the constraint as a test, it changes the stack, until the stack is
empty and an exception is raised.)

On the other hand, this is fine:
\begin{verbatim}
always { e == stack.empty? }
\end{verbatim}

Babelsberg/R will keep track of the dependencies, so that if the stack
is changed by pushing or popping an item, \verb|e| will be kept
up-to-date.

\subsection{Random Number Generator Examples}

A constraint that uses the built-in \verb|rand| method is disallowed
because it calls a system method (Section \ref{sec:system-calls}).
However, it would be easy enough to define our own random number
generator in Ruby:

\begin{verbatim}
class Random
  def init
    @seed = 0
    @a = 1664525
    @c = 1013904223
    @m = 2**32
  end
  def rand
    @seed = (@a*@seed + @c) % @m
    @seed/@m.to_f
  end
end

Then the constraint in this code
r = Random.new
always { x == r.rand.? }
\end{verbatim}

is also disallowed, but not because it uses a random number generator
\emph{per se}, but rather because it uses a method with a side effect.
(If we changed it to use a strange version, with a \verb|rand| method
that just returned the number but didn't update the seed, and a
separate \verb|tick| method that updated the seed, it would be OK.)

\section{Alternatives for the Restrictions}
\label{sec:alternatives}

It seems that the restriction on not creating new objects if the
constraint uses identity constraints is the only one that is
problematic --- the others seem reasonable.  It works, but is stronger
than necessary.

An intuition for a more lenient version of the restriction could be
that if there any new objects are created, such creation must be
``benign'', for example, the new object might be a temporary one used
only in one of the methods called, which then is no longer accessible
after the method returns (and so would be garbage collected).

But what does ``benign'' mean in this context?  One idea is that
identity constraints cannot leak new objects.  But this seems to not
quite work.  Consider:

\begin{verbatim}
def test
  p = Point.new(10,20)
  always {p is? Point.new(10,20)}
  return 42
end

always { test == 42 }
\end{verbatim}

Here, we aren't leaking the new points, but there is still an
unsatisfied constraint in the middle of the \verb|test| method, which
the solver would probably view as satisfiable.

Another approach might be to eliminate the restriction altogether, and
just include a test after the code to satisfy any identity constraint,
to check that in fact the identity constraint holds.  We don't have
any clear use cases for identity constraints that also need to create
new objects, and adding the check has a slight performance penalty ---
on the other hand, it would simplify the language description if it
works.

%% Can we come up with a better example, for instance one in which the
%% return value is different if the identity constraint is correctly
%% solved?

\bibliographystyle{plain}
\bibliography{constraints}

\end{document}
